\subsection*{Implicit Masking}

There are some code structures in Croqtile that implies implicit masking. One typical situation is the multi-level \texttt{parallel-by} block. The below code showcases an example:

\begin{lstlisting}[language=croqtile,caption={Implicit Masking}]
__co__  void foo() {
  parallel p by 6 {
    // parallel-level-0:
    //   code here is as if implicitly masked 'if (q == 0)'
    parallel q by 2 {
      // parallel-level-1:
      //   code here is for all 'q' threads
    }
    // parallel-level-0:
    //   code here is as if implicitly masked 'if (q == 0)'
  }
}
\end{lstlisting}

In this example, there exists two-level of parallelization. The code inside \texttt{parallel p by 6} block but outside \texttt{parallel q by 2} blocks is only for \emph{parallel-level-0}. However, for target like , it could invoke 2 parallel threads in practise to execute either code inside \emph{parallel-level-0} or \emph{parallel-level-1}. Therefore, the \emph{parallel-level-0}-only part is as if guarded with a C++ block \texttt{if (q == 0)}.

Conditional \texttt{foreach} Block Since It is possible to guard the \texttt{foreach} loop with conditions. One direct
